\section{Ngành xếp chuỗi suy luận vào nhóm tự diễn giải được}
\label{sec:ch17-nganh-doc-nguoc-lai}

\index{diễn giải nội tại!chuỗi suy luận bị xếp vào đây|(}%
Mục trước đếm ba lần bài mở đường từ chối tuyên bố về khả năng diễn giải. Mục
này đọc một khảo sát xuất bản ba năm sau đó, và chỗ nó xếp chuỗi suy luận là
đúng chỗ mà ba lần từ chối kia loại trừ.

\index{diễn giải nội tại!ba nhóm của khảo sát}%
Khảo sát ấy hỏi mô hình ngôn ngữ giúp được gì cho việc giải thích các mô hình AI,
và nó chia các kỹ thuật thành ba nhóm:
\tn{giải thích hậu nghiệm}{post-hoc explanation}, diễn giải nội tại, và giải
thích lấy con người làm trung tâm~\autocite{p16llmxai}.
\tn{Diễn giải nội tại}{intrinsic interpretability} là nhóm mà hình 2 của bài mô
tả bằng một dòng, thiết kế mô hình ngôn ngữ sao cho nó tự diễn giải được, và cho
hai ví dụ: kiến trúc mô hình trong suốt, và khả năng diễn giải dựa trên
\tn{cơ chế chú ý}{attention mechanism}. Chuỗi suy luận không có mặt trong hình.
Nó vào nhóm ấy ở phần thân, trang 4~\autocite{p16llmxai}:

\begin{quote}
Một kỹ thuật đáng chú ý trong nhóm này là suy luận theo chuỗi suy luận, kỹ thuật
chia các bài toán phức tạp thành những bước nhỏ hơn, dễ quản lý hơn, cho phép mô
hình đi theo một tiến trình logic tới kết luận.%
\footnote{Nguyên văn: \enquote{A notable technique in this category is Chain of
Thought (CoT) reasoning [38], which breaks down complex problems into smaller,
manageable steps, allowing models to follow a logical progression toward
conclusions.}}
\end{quote}

Nghĩa của nước đi ấy là gì thì phải đọc theo định nghĩa của chính nhóm. Một
phương pháp hậu nghiệm được kiểm bằng cách hỏi nó có phản ánh đúng mô hình
không; đó là toàn bộ Phần~IV. Một mô hình \emph{tự} diễn giải được thì không cần
phép kiểm ấy, vì cấu trúc của nó đã là lời giải thích. Xếp chuỗi suy luận vào
nhóm thứ hai, do đó, không phải một cách sắp xếp cho gọn danh mục. Nó là lời
tuyên bố rằng ở đây không có gì phải đối chiếu.

\index{diễn giải nội tại!khảo sát dẫn ai cho nước đi ấy}%
Đáng chú ý là bài dẫn ai cho câu vừa trích. Số hiệu trong ngoặc không trỏ về bài
mở đường mà trỏ về một khảo sát khác về chuỗi suy luận, của Chu và cộng
sự~\autocite{p16llmxai}. Bài mở đường \emph{có} nằm trong danh mục tài liệu của
khảo sát, dưới đúng nhan đề của nó, nên đây không phải chuyện không biết tới.
Câu phụ lục C.2 mà mục~\ref{sec:ch17-bai-mo-duong} trích nằm trong một bài mà
khảo sát này có trên kệ và không mở tới.

\index{diễn giải nội tại!bảng 1 dẫn bài đo ngược lại}%
Chỗ hỏng thật thì nằm trong một bảng. Bảng 1 của khảo sát, trang 6, tóm tắt cả
bài thành mười dòng, mỗi dòng một chủ đề, một cột ý chính và một cột nguồn. Dòng
\enquote{Intrinsic interpretability} có cột ý chính ghi \enquote{Chain of thought
reasoning, guided templates for logical progression} và cột nguồn liệt kê sáu số
hiệu~\autocite{p16llmxai}. Một trong sáu số ấy trỏ về bài của Lanham và cộng sự,
nhan đề \emph{Measuring Faithfulness in Chain-of-Thought Reasoning}. Rà toàn văn
bài khảo sát, số hiệu ấy xuất hiện đúng hai lần: một lần trong ô bảng vừa nói, và
một lần trong chính mục danh mục tài liệu của nó. Nó không xuất hiện lần nào
trong phần thân.

Nghĩa là bài duy nhất trên kệ của khảo sát từng đo xem chuỗi suy luận có trung
thực hay không được dẫn đúng một lần, để chống đỡ cho lời khẳng định rằng chuỗi
suy luận làm mô hình tự diễn giải được, và kết quả nó đo thì không được nhắc.

\index{diễn giải nội tại!bài Lanham đo được gì}%
Sách đọc bài ấy, vì không đọc thì không phán được câu vừa rồi. Đây là lần thứ tư
sách phải mở một bài ngoài danh mục 32 bài, sau bảy khóa của
chương~\ref{ch:do-do-trung-thuc}, một khóa của
chương~\ref{ch:giai-thich-duoi-tan-cong} và một khóa của
chương~\ref{ch:nut-that-khai-niem}, và lý do lần này là lý do của lần thứ ba: một
bài trong danh mục nói một điều \emph{về} một bài khác, và chương phải xử câu ấy.
Bài của Lanham và cộng sự can thiệp vào chuỗi suy luận, chẳng hạn bằng cách chèn
lỗi vào giữa chuỗi hoặc viết lại chuỗi bằng lời khác, rồi xem dự đoán của mô hình
đổi thế nào. Kết quả đo của nó nằm trong một câu của phần tóm
tắt~\autocite{mcotfaith}:

\begin{quote}
Khi các mô hình trở nên lớn hơn và có năng lực hơn, chúng sinh ra suy luận
\emph{kém trung thực hơn} trên hầu hết các nhiệm vụ chúng tôi khảo sát.%
\footnote{Nguyên văn: \enquote{As models become larger and more capable, they
produce less faithful reasoning on most tasks we study.}}
\end{quote}

Câu ấy một mình thì dễ bị đọc quá tay, nên phải đọc kèm câu kết luận của cùng
phần tóm tắt, và sách in cả hai vì mỗi câu chặn một cách đọc sai khác nhau:

\begin{quote}
Nhìn chung, kết quả của chúng tôi gợi ra rằng chuỗi suy luận \emph{có thể} trung
thực nếu các điều kiện như quy mô mô hình và nhiệm vụ được chọn cẩn thận.%
\footnote{Nguyên văn: \enquote{Overall, our results suggest that CoT can be
faithful if the circumstances such as the model size and task are carefully
chosen.}}
\end{quote}

\index{diễn giải nội tại!chuỗi suy luận bị xếp vào đây|)}%
Câu thứ hai giữ cho câu thứ nhất khỏi bị đọc quá tay, và sách in nó vì đúng lý do
ấy. Bài không kết luận rằng chuỗi suy luận không trung thực; nó kết luận rằng độ
trung thực của chuỗi phụ thuộc vào mô hình nào và nhiệm vụ nào. Nhưng phụ thuộc
vào mô hình chính là điều một kỹ thuật \emph{tự} diễn giải được không được phép
có: nếu tính chất ấy đổi theo quy mô, và đổi theo chiều xấu đi, thì nó không phải
tính chất của cách nhắc mà là tính chất của từng cặp mô hình và nhiệm vụ, và cần
đo ở từng cặp. Đó là lời phát biểu hẹp và đúng, và nó đủ để bác dòng bảng: bài
được dẫn ở đó không ủng hộ cái nó được dẫn để ủng hộ.
